\section{Introduction}

A multimodal language model can look at a scene, plan a motion, and call a tool
to execute it. But producing a plausible plan, or declaring the task finished,
establishes nothing about the resulting physical state. Evaluating these models
as robot controllers therefore means scoring the environment rather than the
agent's report, and accounting for the inference the attempt consumed.

We vary two factors and fix everything else: which model drives the robot, and
how much reasoning it is asked to spend per decision. Reasoning effort is sold
as a quality dial, but whether it buys physical competence, and at what price,
is unmeasured. Our protocol fixes the observation interface, the manipulation
prompt, the simulator horizon, and a per-episode allowance of model calls whose
cost is recorded; success is LIBERO's own predicate, never the agent's claim.
We evaluate on LIBERO-10~\citep{liu2023libero} through the Inspect Robots
agent~\citep{inspectrobots} and a custom embodiment adapter, with the model
selecting end-effector pose targets directly and no task-specific policy
trained (Figure~\ref{fig:interface}).

Recent evaluations of this model generation take the opposite cut.
\citet{galbot2026astra} report a gap between competent task-level decisions and
reliable control, with GPT-6 Astra at 26\% on RoboDojo when it emits control
targets itself against 48\% when it corrects a learned policy;
\citet{roboicl2026} and \citet{vlmagents2026} hold one model fixed and engineer
its context. None varies the model, and none varies reasoning effort. We take
the hardest corner --- one frozen model, no learned policy, no demonstrations,
one shared interface --- because it isolates what a model contributes alone.
Trained vision-language-action policies exceed 90\% on LIBERO-10, so these
rates measure a different thing, not a worse attempt at the same one.

Across fifteen configurations --- five models $\times$ three reasoning levels
--- and 1,500 episodes, we find: the best configuration reaches 44\% and the
campaign averages 19.0\%; reasoning effort raises success for three models but
not Claude Opus 5; accuracy and cost are near-independent, with
two configurations tied at 33\% differing threefold in cost per success; and
failure is uneven, with one model never succeeding in 300 episodes and three of
ten tasks solved by no configuration.
